\setchapter{2}{Binary Choice Models}
\ChapterHeading

\section{Linear probability model}
In many applications, the variate of interest is binary, i.e., takes values 0 and 1.
Let's denote $X = (X_1,\ldots,X_k)'$ and the objective of the regression model is to estimate $E[Y|X]$
where $E[Y|X] = P(Y=1|X)$.\\

In the \Term{linear probability model} we assume that
\begin{equation}
    P(Y=1|X) = \beta_1 + \beta_2 X_2 + \cdots + \beta_k X_k
\end{equation}
thus the \Emph{interpreation of $\beta_j$} is the \KeyIdea{change in the probability of success when $x_j$ changes}.
\begin{equation}
    \frac{\partial P(Y=1|X)}{\partial X_j} = \beta_j,\quad j=2,\ldots,k
\end{equation}
and the \Emph{predicted $Y$} is the \KeyIdea{predicted probability of success}.\\

The linear probability model is estimated using \Emph{OLS}.
But there is the potential problem that the fitted values can be outside $[0,1]$.
Even without predictions outside $[0,1]$, we may estimate effects that imply a change in $x$ changes the probability by more than $+1$ or $-1$.
\Info{$\hat{\beta}$ too big that would not make sense.}
Also, this model \Emph{violates the assumption of homoskedasticity}, 
therefore we should use the Eicker-Huber-White \Emph{robust standard errors} to make inference.
\Info{Homoskedasticity states that $\Var[Y|X]$ does not vary with $x$ but here it necessarity depends on $x$.}
\begin{proof}
    \begin{equation*}
        \Var[Y|X] = \underbrace{E[Y^2|X]}_{(1)} - E[Y|X]^2
    \end{equation*}
    (1):
    \begin{equation*}
        E[Y^2|X] = 1^2\times P(Y=1|X)+0^2\times P(Y=0|X) = P(Y=1|X)
    \end{equation*}
    Now we have,
    \begin{equation*}
        \Var[Y|X] = P(Y=1|X) - P(Y=1|X)^2 = \underbrace{P(Y=1|X)}_{\beta_1+\beta_2 x_2+\cdots+\beta_k x_k}[1-P(Y=1|X)]
    \end{equation*}
    So we see that the variance depends on the regressors.
\end{proof}

\section{Index models for binary response}
The alternative to the linear model is to assume that $E[Y|X] = P(Y=1|X) = G(X'\beta_0)$,
where $G(\cdot)$ is bounded between 0 and 1. Thus, we have,
\begin{equation}
    Y = \begin{cases}
        1 & \text{with probability }G(X'\beta_0)\\
        0 & \text{with probability }1-G(X'\beta_0)
    \end{cases}
\end{equation}
with the log-likelihood function of
\begin{equation}
    \log\{L(\beta)\} = \sum_{i=1}^{n} Y_i\log(G(X_i'\beta)) + \sum_{i=0}^{n}(1-Y_i)\log(1-G(X_i'\beta))
\end{equation}
The \Emph{MLE} should be used to find the estimators, $\hat{\beta}_{ML}$.\\

Note that it is a system of non-linear equations to be solved to find $\hat{\beta}_{ML}$
thus we have to resort to numerical methods. \Info{There is no closed form solution for this estimator.}
\Emph{Consistency and asymptotic normality} follow the general results described for MLE under regulairty conditions.
Essentially, the \Emph{main requirement for consistency} is $E[Y|X] = P(Y=1|X) = G(X'\beta_0)$.\\

In a correctly specified model, $\hat{\beta}$ is consistent and asymptotically normally distributed
with variance-covariance matrix $[\mcI(\beta_0)]^{-1}$,
\begin{equation*}
    \sqrt{n}(\hat{\beta}-\beta_0)\xrightarrow{d}\mcN(0,[\mcI(\beta_0)]^{-1})
\end{equation*}
\Emph{Inference} is done using the Wald, Likelihood Ratio, and Lagrange multiplier statistics as usual.

\subsection{$G(X'\beta_0)$ Choices}
The most popular choices are
\begin{itemize}
    \item The \Term{Logit Model}
    \begin{equation}
        G(X'\beta_0) = \frac{\exp(X'\beta_0)}{1+\exp(X'\beta_0)}
    \end{equation}
    \item The \Term{Probit Model}
    \begin{equation}
        G(X'\beta_0) = \Phi(X'\beta_0)
    \end{equation}
    where $\Phi(z)=\frac{1}{\sqrt{2\pi}}\int_{-\infty}^{z}\exp(-\frac{u^2}{2})\,du$ is the Standard Normal Distribution function.
\end{itemize}
Both models are non-linear and require \Emph{MLE}. Other possible models are
\begin{itemize}
    \item The \Term{Log-Weibull} Distribution
    \begin{equation*}
        G(X'\beta_0) = \exp(-\exp(X'\beta_0))
    \end{equation*}
    \item The \Term{Complementary log-log Model}, Gompertz Distribution
    \begin{equation*}
        G(X'\beta_0) = 1- \exp(-\exp(X'\beta_0))
    \end{equation*}
    \item \begin{equation*}
        G(X'\beta_0) = \Phi(X'\beta_0)^{\tau}
        \end{equation*}
    where $\tau$ needs to be estimated.
    \item \begin{equation*}
        G(X'\beta_0) = 1- \exp(-\omega\exp(X'\beta_0))^{-\frac{1}{\omega}}
    \end{equation*}
    where $\omega$ needs to be estimated. 
    Note that for $\omega = 1$ we have the Logit model, 
    and $\lim_{\omega\to 0}G(X'\beta_0) = 1- \exp(-\exp(X'\beta_0))$.
\end{itemize}
\Info{Estimating additional parameters of the distribution allows the data to tell us the true model
which is a major benefit.}\\

In statistics, a more common interpretation of the coefficients 
is in terms of \Term{marginal effects} on the odds ratio rather than on the probability.
$\frac{P(Y=1|X)}{(1-P(Y=1|X))}$ measures the \KeyIdea{probability of yes relative to no}
and is called the \Emph{odds ratio} or relative risk.
For the logit model, the log-odds ratio is linear in the regressors.
\begin{proof}
    \begin{align*}
        P(Y=1|X) & = \frac{\exp(X'\beta_0)}{1+\exp(X'\beta_0)} \Rightarrow\\
        \frac{P(Y=1|X)}{1-P(Y=1|X)} & = \frac{\frac{\exp(X'\beta_0)}{1+\exp(X'\beta_0)}}{1-\frac{\exp(X'\beta_0)}{1+\exp(X'\beta_0)}} = \exp(X'\beta_0)\Rightarrow\\
        \log\left(\frac{P(Y=1|X)}{1-P(Y=1|X)}\right) & = X'\beta_0
    \end{align*}
\end{proof}

\subsection{Latent Variable Threshold (LVT) Model}
Define a latent variable
\begin{equation}
    Y^* = X'\beta_0 + \epsilon_t
\end{equation}
where $Y^*$ is unobserved, thus a \Term{latent variable}.
Further assume that $\epsilon$ independent of $X$, $E[\epsilon]=0$ and $\Var(\epsilon)=\sigma^2$ with a distribution function of $F(\cdot)$.
The observation of $Y$ is given by
\begin{equation}
    Y =
    \begin{cases}
        1 & \text{if }Y^* > \lambda\\
        0 & \text{if }Y^* \leq\lambda
    \end{cases}
\end{equation}
where $\lambda$ is the threshold.\\

The probability of choosing the option is thus given by,
\begin{align*}
    P(Y=1|X) & = P(Y^*>\lambda|X) = P(X'\beta_0+\epsilon>\lambda|X)\\
    & = P(\epsilon > -X'\beta_0+\lambda|X)\\
    & = 1 - P(\epsilon \leq -X'\beta_0+\lambda|X)\\
    & \text{independence between }\epsilon, X\\
    & = 1 - \underbrace{F(-X'\beta_0+\lambda)}_{\text{cumulative distribution function}}\\
    & = G(X'\beta_0)
\end{align*}
with $G(z) = 1 - F(-z+\lambda)$.\\

\Term{First identification problem}: 
\begin{equation}
    P(Y=1|X) = 1 - F(-X'\beta_0+\lambda)
\end{equation}
Notice that if $X_1 = 1$ meaning that \Emph{there is an intercept in the model}, 
it is not possible to estimate $-\beta_1$ and $\lambda$ at the intercept separately. 
Solution is to \KeyIdea{set $\lambda = 0$}. If $\lambda = 0$, and $\epsilon$ has a symmetric distribution around 0,
then $G(z) = F(z)$ as in this case $1-F(-z) = F(z)$.\\

\Term{Second identification problem}: 
\begin{equation}
    \frac{Y^*}{a} = X'\beta_0^*+\frac{\epsilon}{a}
\end{equation}
If we divide $Y^*$ from (24) by a constant $a$ and $a>0$, 
notice that the definition of the observable variable $Y$ (25) does not change.
\Info{Changing the scale $\to$ changing of variance but the $Y$ does not change $\Rightarrow$
we cannot identify the scale of $Y^*$.} This implies that \Emph{we cannot identify the variance of $\epsilon$}.
\Info{Error term variance not observed and cannot be identified and estimated.}
For given $\beta_0^*$, value of $\beta_0$ depends on $a$, meaning $\beta_0$ is identified up to a scale factor.
The solution is to \KeyIdea{normalize distribution of $\epsilon$} by fixing $\sigma^2$ at a given number $\Rightarrow$ assume $\sigma^2$ is known.
\Info{Normally set to 1 in practice.}

\subsection{Random Utility Models}
Suppose an individual has to choose between alternatives $a$ and $b$ with utilities $U^a$ and $U^b$ which are unobservable.
But the characteristics of the observations are observable thus we have,
\begin{align*}
    U^a & = X'\beta_a + u_a\\
    U^b & = X'\beta_b + u_b
\end{align*}
We also observe the chosen alternative, say $a$, which we indicate this choice as $Y=1$.
Then we know that,
\begin{align*}
    P(Y=1|X) & = P(U^a > U^b|X) = P(X'\beta_a+u_a > X'\beta_b+u_b)\\
    & = P(u_a-u_b>X'(\beta_b-\beta_a)|X)\\
    & = P(\epsilon > -X'\beta_0|X) = 1 - F(-X'\beta_0)
\end{align*}
with $\epsilon = u_a-u_b$ and $\beta_0=\beta_a-\beta_b$.

\section{Marginal effects and average partial effects}

It is important to note that the parameters of the model, like those of any nonlinear regression model,
are not necessarily the marginal effects. In general,  the partial effect associated with variable $k$, $X_{ik}$ is given by,
\begin{equation}
    \Delta E[Y_i|X_i] \equiv \Delta P(Y_i = 1|X_i) \simeq \frac{\partial E[Y_i|X_i]}{\partial X_{ik}}\Delta X_{ik} = \frac{\partial G(X'_i\beta_0)}{\partial X_{ik}}\Delta X_{ik}
\end{equation}
Let $\omega = X'_i\beta_0 = \beta_1X_{i2}+\cdots+\beta_k X_{iK}$, then,
\begin{equation}
    \frac{\partial E[Y_i|X_i]}{\partial X_{ik}} = \frac{\partial G(\omega)}{\partial(\omega)}\frac{\partial(\omega)}{\partial X_{ik}}
    = g(\omega)\beta_k = g(X'_i\beta_0)\beta_k
\end{equation}
where $g(\omega) = \frac{\partial G(\omega)}{\partial\omega}$. Thus, we need to evaluate
\begin{equation}
    \Delta e[Y_i | X'_i\beta_{0i}] \equiv \Delta P(Y_I=1|X_i) \simeq g(X'_i\beta_0)\beta_k\Delta X_{ik}
\end{equation}

Now, to \Emph{compute marginal effects} of variable $X_k$ assuming that \Emph{$X_{ik}$ is continuous}, we have options 
\begin{itemize}
    \item Marginal effect at a \KeyIdea{particular value $X_i = X_0$}
    \begin{equation*}
        g(X'_0\hat{\beta}_{ML})\hat{\beta}_{ML,i}\Delta X_{ik}
    \end{equation*}
    The most common choices of $X_0$ are the vector of sample means. This does not make sense if some of the regressors are dummy variables.
    Or the vector of sample medians.
    \item Alternative is to compute \KeyIdea{Average Partial Effects} which is the sampel average of the individual marginal effects,
    \begin{equation*}
        \frac{1}{n}\sum_{i=1}^{n}g(X_i'\hat{\beta}_{ML})\hat{\beta}_{ML,k}\Delta X_{ik}
    \end{equation*}
\end{itemize}
The options to compute the marginal effects on a \Emph{dummy variable} $d$, we have the same options 
\begin{itemize}
    \item Marginal effect at a \KeyIdea{particular value $X_i = X_0$}
    \begin{equation*}
        P(\widehat{Y_i=1|X_0,d_i=1}) - \widehat{P(Y_i=1|X_0,d_i=0)}
    \end{equation*}
    Note that here $X_0$ denotes all variables in the model at the particular value $X_0$ except $d$.
    \item Average partial effect is thus given by 
    \begin{equation*}
        \frac{1}{n}\sum_{i=1}^{n}P(\widehat{Y_i=1|X_i,d_i=1}) - \frac{1}{n}\sum_{i=1}^{n}\widehat{P(Y_i=1|X_i,d_i=0)}
    \end{equation*}
    Note that here $X_i$ denotes all variables in the model except $d$.
\end{itemize}

\section{Simple specification tests}
\subsection{Misspecification}
If $G(\cdot)$ is \Emph{misspecified}, then $\hat{\beta}_{ML}$ is \Emph{inconsistent}. 
Perform \KeyIdea{RESET-type test} by testing \Emph{$H_0: \delta_1 = \delta_2 = 0$} in the model 
\begin{equation*}
    E[Y_i|X_i] = G(X'_i\beta_0+\delta_1(X'_i\hat{\beta}_{ML})^2+\delta_2(X'_i\hat{\beta}_{ML})^3),\quad i=1,\ldots,n
\end{equation*}
\Info{This is a normality test in the probit if $\epsilon_i\sim\mcN(0,1)\to f(\cdot)=\Phi(\cdot)$.
$E[Y_i|X_i]=X'_i\beta$ estimated by OLS and the estimator $\hat{\beta}$ and 
the model $Y_i = X'_i\beta + \delta_0(X'_i\beta)^2+\delta_1(X'_i\beta)^3$ and test $H_0: \delta_0 = 0$ and $\delta_1=0$, if we reject $H_0 \to$ model is misspecified.}\\

The RESET test can be tested against mroe general parametric specifications which include additional shape parameters.
Consider $G(\cdot) = \Phi(\cdot)^\tau$, use the score statistic to test $H_0: \tau=1$, Probit. Or, consider $G(\cdot) = 1-(1+\omega\exp(\cdot))^{-\frac{1}{\omega}}$,
use the score statistic to test $H_0: \omega=1$, Logit.

\subsection{Heteroskedasticity}
In the LVT model, heteroskedasticity leads to misspecification of the conditional mean of $Y$.
\begin{proof}
    Proof of misspecification for LVT.\\
    The goal is to show that $E[Y|X] \neq G(X'\beta_0)$. As previously shown, $G(a) = 1 - F(a)$.
    \begin{align*}
        P(Y=1|X) & = P(Y^* > 0|X) = P(X'\beta_0+k\,h(Z'\gamma_0)\epsilon > 0|X)\\
        & = P\left(\epsilon > \frac{-X'\beta_0}{k\,h(Z'\gamma_0)}|X\right)\\
        & \text{notice $X'$ and condition on $X$, we can drop the condition}\\
        & = P\left(\epsilon > \frac{-X'\beta_0}{k\,h(Z'\gamma_0)}\right)\\
        & = 1 - P\left(\epsilon \leq \frac{-X'\beta_0}{k\,h(Z'\gamma_0)}\right)\\
        & = 1 - F\left(\epsilon \leq \frac{-X'\beta_0}{k\,h(Z'\gamma_0)}\right)\\
        & = G\left(\epsilon \leq \frac{-X'\beta_0}{k\,h(Z'\gamma_0)}\right) \neq G(X'\beta_0)\to \text{misspecified}
    \end{align*}
\end{proof}
Consider a latent variable, 
\begin{equation}
    Y^* = X'\beta_0 + k\times h(Z'\gamma_0)\epsilon
\end{equation}
under the assumptions of $\epsilon$ independent of $X$, $E[\epsilon]=0$ and $\Var(\epsilon)=1$ and distirbution function $F(\cdot)$, 
$Z$ are a vector function of $X$ of size $d$ and $h$ any function with $h>0, h(0) = 1, h'(0)\neq0$.
\begin{proof}
    Proof of heteroskedasticity.\\
    The idea is to prove that the variance of the error term is not constant but a function of $X$.
    \begin{gather*}
        u = k\times h(Z'\gamma_0)\epsilon\\
        e[u|X] = E[k\times h(Z'\gamma_0)\epsilon|X] = k\times h(Z'\gamma_0) \underbrace{E[\epsilon|X]}_{=0} = 0\\
        \Var(u|X) = \Var(k\times h(Z'\gamma_0)\epsilon|X) = k^2[h(Z'\gamma_0)]^2\underbrace{\Var(\epsilon|X)}_{=1} = k^2[h(Z'\gamma_0)]^2
        \to \text{Heteroskedasticity}
    \end{gather*}
\end{proof}

We can construct a \KeyIdea{LM test} based on the generalized residuals to test $H_0: \gamma_0=0$ which implies \Emph{homoskedasticity}.
The LM test statistic is calculated as 
\begin{equation*}
    \xi_{LM} = \iota'S(S'S)^{-1}S'\iota\sim\chi^2(d)
\end{equation*}
This is asymptotically equivalent to a \KeyIdea{Wald test} with $H_0: \gamma_0=0$ in the model 
\begin{equation*}
    E[Y_i|X_i] = G(X'\beta_0 + (X'_i\hat{\beta}_{ML})Z'_i\gamma_0),\quad i=1,\ldots,n
\end{equation*}

\subsection{Goodness of fit}
Unlike the linear probaility model, we cannot compute an $R^2$ to judge the goodness of fit.
One possibility is to use the \KeyIdea{pseudo $R^2$} based on the log likelihood. We can also look at the percent correctly predicted.
\Info{We can just use the percentage of correct predictions and we can adjust the threshold. This is more for machine learning when prediction matters.
For econometrics, the interpretation of $\beta$ is more important, thus the specification test is more important.}


% \begin{definition}
% Theorem
% Lemma
% \begin{Example}
% \begin{proof}
%\newcommand{\KeyIdea}  % 关键想法：浅绿底
%\newcommand{\Term}     % 专业名词：品牌深绿字
%\newcommand{\Emph}     % 强调词汇：红
%\newcommand{\Confuse}    % 不懂/存疑处
%\newcommand{\Info}     % 说明或注解